\documentclass[12pt]{article}
\usepackage[utf8]{inputenc}
\usepackage[T1]{fontenc}
\usepackage{amsmath,amssymb,mathtools}
\usepackage{geometry}
\usepackage{booktabs}
\usepackage{xcolor}
\usepackage[colorlinks]{hyperref}
\hypersetup{citecolor=blue,linkcolor=blue}
\allowdisplaybreaks
\newcommand{\CV}[1]{{\color{blue}#1}}
\newcommand{\todo}[1]{{\color{red}#1}}

\title{Documentation of the model implemented in Python}
\author{Comparison with the paper formulation}
\date{}

\begin{document}
\maketitle

\section{Introduction}

This document describes the MIQP model implemented in the file \texttt{model.py}
and highlights the main differences with respect to the mathematical formulation
in the paper \emph{Drone Coverage Path Planning} (Amorosi, Dell'Olmo, Puerto, Valverde).

\subsection*{Common notation}

\begin{itemize}
  \item $\mathcal R$: set of regions.
  \item $\mathcal V'$: vertices (interruption points in each region).
  \item $\mathcal E^{int} = \mathcal E_{RL}^{int} \cup \mathcal E_{LR}^{int}$: intra-region edges (RL = forward, LR = backward).
  \item $\mathcal E^{ext}$: inter-region edges.
  \item $\mathcal T_r$: polygonal chains available for region $r$.
  \item $\mathcal S_t$: segments of chain $t$.
  \item $Q^s, Q^{s+1}$: endpoints of segment $s$.
  \item $L(s) = \|Q^{s+1} - Q^s\|$: length of segment $s$.
  \item $E^{xy} = \frac12 c A$, $E^{z} = \frac12 c A$: drag constants.
  \item $\rho_t$: air density at the height of chain $t$.
  \item $\nu_{t,s}^d$: actual drone speed over segment $s$ of chain $t$.
  \item $W^{max}$: maximum drone endurance.
\end{itemize}

\newpage
\section{Variables of the implemented model}

\subsection{Binary variables}
\begin{center}
\begin{tabular}{lll}
\toprule
Variable & Domain & Description \\
\midrule
$\mu_v^{ts}$ & $\mathcal V' \times \{(t,s): t\in\mathcal T_r, s\in\mathcal S_t\}$ & 1 if vertex $v=(r,i)$ is assigned to segment $s$ of chain $t$. \\
$\alpha_t$ & $\{t\in\mathcal T_r : r\in\mathcal R\}$ & 1 if chain $t$ is selected for its region. \\
$x_{uv}^o$ & $\mathcal E \times \mathcal O$ & 1 if the drone traverses $(u,v)$ in operation $o$. \\
$y_v^o$ & $\mathcal V \times \mathcal O$ & 1 if vertex $v$ is visited in operation $o$. \\
$\zeta^o$ & $\mathcal O$ & 1 if all vertices have been visited before operation $o$. \\
$\varphi_{rr'}$ & $\mathcal R \times \mathcal R$, $r\neq r'$ & 1 if the drone moves from $r$ to $r'$ at the height of $r$. \\
$\psi_{rr'}$ & $\mathcal R \times \mathcal R$, $r\neq r'$ & 1 if the height of $r$ is lower than that of $r'$. \\
\bottomrule
\end{tabular}
\end{center}

\subsection{Continuous and integer variables}
\begin{center}
\begin{tabular}{lll}
\toprule
Variable & Range & Description \\
\midrule
$P_v$ & $\mathbb R^3$ & Coordinates of the point associated with vertex $v$. \\
$\gamma_v$ & $[0,1]$ & Position within the selected segment. \\
$\lambda_v$ & $[0, S_t]$ & Global parametrization on the selected chain. \\
$dfs_v$ & $[0, L_{max}]$ & Distance from the start of the chain. \\
$efs_v$ & $[0, L_{max}\cdot 10]$ & Energy from the start of the chain. \\
$d_{uv}^{xy}$ & $[0, D_{max}]$ & Inter-region horizontal distance. \\
$d_{uv}^z$ & $[0, H_{max}]$ & Inter-region vertical distance. \\
$d_{uv}$ & $[0, D_{max}+H_{max}+L_{max}]$ & Total distance between $u$ and $v$. \\
$w_{uv}$ & $[0, \infty)$ & Energy between $u$ and $v$. \\
$k^o$ & $\{0,\dots,|\mathcal V'|\}$ & Number of vertices visited in operation $o$. \\
$u_v^o$ & $\{1,\dots,|\mathcal V'|\}$ & MTZ variable for subtour elimination. \\
$\rho_r$ & $[0,\rho(0)]$ & Selected density for region $r$. \\
$\eta_v^{ts}$ & $[0,1]$ & Linearization of $\gamma_v\cdot\mu_v^{ts}$. \\
% Various linearization auxiliaries
\bottomrule
\end{tabular}
\end{center}

\section{Location constraints}

They match the paper in structure. The main difference is in the fixing
of the START and END vertices:

\begin{align}
  P_v &= \sum_{t\in\mathcal T_r}\sum_{s\in\mathcal S_t}
        \bigl[ \mu_v^{ts} Q^s + \eta_v^{ts} (Q^{s+1}-Q^s) \bigr],
        && \forall v=(r,i)\in\mathcal V', \tag{LC1}\\
  P_v^z &= \sum_{t\in\mathcal T_r}\sum_{s\in\mathcal S_t} \mu_v^{ts} \cdot (\text{height of }t),
        && \forall v=(r,i)\in\mathcal V', \tag{LC1z}\\
  \sum_{s\in\mathcal S_t} \mu_v^{ts} &= \alpha_t,
        && \forall v=(r,i)\in\mathcal V',\; \forall t\in\mathcal T_r, \tag{LC2}\\
  \sum_{t\in\mathcal T_r} \alpha_t &= 1,
        && \forall r\in\mathcal R, \tag{LC3}\\
  \lambda_v &= \sum_{t\in\mathcal T_r}\sum_{s\in\mathcal S_t} (s-1)\mu_v^{ts} + \gamma_v,
        && \forall v\in\mathcal V', \tag{LC4}\\
  \lambda_v &= 0,
        &&\forall v=(r,1),(r,2), \tag{LC5}\\
  \mu_v^{t0} &= \alpha_t,
        &&\forall v=(r,1),(r,2),\; \forall t\in\mathcal T_r, \tag{LC5b (new)}\\
  \lambda_u &\le \lambda_v,
        &&\forall (u,v)\in\mathcal E_{RL}^{int}, \tag{LC6}\\
  \lambda_u &= \lambda_v,
        &&\forall (u,v)\in\mathcal E_{LR}^{int}, \tag{LC7}\\
  \lambda_v &= \sum_{t\in\mathcal T_r} S_t \alpha_t,
        &&\forall v=(r,2R+1),(r,2R+2), \tag{LC8}\\
  \mu_v^{t,S_t-1} &= \alpha_t,
        &&\forall v=(r,2R+1),(r,2R+2),\; \forall t\in\mathcal T_r. \tag{LC8b (new)}
\end{align}

Constraints LC5b and LC8b explicitly fix the segment (first or last)
of the START and END vertices, a simplification with respect to the paper
where this fixing was done implicitly through $\lambda_v$ and the domain
of $\mu$.

\section{Drone path constraints}

They match the paper except for subtour elimination:

\begin{align}
  \sum_{v\neq v_0} x_{v_0 v}^o &= 1 - \zeta^o,
        && \forall o\in\mathcal O, \tag{DP1}\\
  \sum_{u} x_{uv}^o - \sum_{u} x_{vu}^o &= 0,
        && \forall v\in\mathcal V',\; \forall o\in\mathcal O, \tag{DP2}\\
  \sum_{v\neq v_0} x_{v v_0}^o &= 1 - \zeta^o,
        && \forall o\in\mathcal O, \tag{DP3}\\
  \sum_{u} x_{vu}^o &= y_v^o,
        && \forall v\in\mathcal V,\; \forall o\in\mathcal O, \tag{DP4}\\
  \sum_{u} x_{uv}^o &= y_v^o,
        && \forall v\in\mathcal V,\; \forall o\in\mathcal O, \tag{DP5}\\
  \sum_{o} y_v^o &= 1,
        && \forall v\in\mathcal V', \tag{DP6}\\
  u_u^o - u_v^o + 1 &\le (|\mathcal V'|-1)(1-x_{uv}^o),
        && \forall u\neq v\in\mathcal V',\; \forall o\in\mathcal O, \tag{DP7 (MTZ)}\\
  \sum_{o} x_{uv}^o &= 1,
        && \forall (u,v)\in\mathcal E_{RL}^{int}. \tag{DP8}
\end{align}

\noindent The paper's main formulation uses the same MTZ constraints. The paper also
describes an alternative edge-based formulation that replaces MTZ with
Dantzig--Fulkerson--Johnson (DFJ) exponential subtour elimination constraints,
separated dynamically via a lazy constraint callback during the branch-and-cut
search (see Section~\ref{sec:edge-model} below). Both are implemented in the
codebase~\footnote{The vertex-based model is \texttt{RingsModel} in
\texttt{mip\_rings.py}, the edge-based model is \texttt{EdgesModel} in
\texttt{mip\_edges.py}.}.

\section{Intra-region distances}

\subsection*{Paper approach (with bilinear $\beta$)}
The paper defines $d_{vv'}^{ss'}$ using bilinear products
$\beta_{vv'}^{ss'} = \mu_v^s \mu_{v'}^{s'}$, resulting in expressions
depending on 4 indices (two vertices, two segments).

\subsection*{Implemented approach (DFS)}
Instead of $\beta$, a continuous variable $dfs_v$ is introduced to represent
the cumulative distance from the start of the chain to vertex $v$:

\begin{align}
  dfs_v &= \sum_{t\in\mathcal T_r}\sum_{s\in\mathcal S_t}
          \bigl[ \text{cum\_before}(s) \cdot \mu_v^{ts}
          + L(s) \cdot \eta_v^{ts} \bigr],
          && \forall v\in\mathcal V', \tag{DFS}\\
  d_{uv} &= dfs_v - dfs_u,
          && \forall (u,v)\in\mathcal E_{RL}^{int}, \tag{IDIST}\\
  d_{uv} &= 0,
          && \forall (u,v)\in\mathcal E_{LR}^{int}. \tag{IDIST0}
\end{align}

where $\text{cum\_before}(s) = \sum_{s' < s} L(s')$, and $\eta_v^{ts} = \gamma_v \cdot \mu_v^{ts}$
is linearized via standard Fortet constraints (binary $\times$ continuous).

This formulation is more compact: it avoids the double products $\mu_v^s\mu_{v'}^{s'}$
and reduces the number of variables and quadratic constraints.

\section{Intra-region energy}

Analogously to distance, a variable $efs_v$ is used for cumulative energy:

\begin{align}
  efs_v &= \sum_{t\in\mathcal T_r}\sum_{s\in\mathcal S_t}
          \bigl[ C_{ts} \cdot \text{cum\_before}(s) \cdot \mu_v^{ts}
          + C_{ts} \cdot \eta_v^{ts} \bigr],
          && \forall v\in\mathcal V', \tag{EFS}\\
  w_{uv} &= efs_v - efs_u,
          && \forall (u,v)\in\mathcal E_{RL}^{int}, \tag{IENER}\\
  w_{uv} &= 0,
          && \forall (u,v)\in\mathcal E_{LR}^{int}, \tag{IENER0}
\end{align}
where $C_{ts} = E^{xy} \rho_t \nu_{ts}^d L(s)$ is the energy per unit
distance on segment $s$ of chain $t$, computed as a constant
(pre-process).

\section{Inter-region distances and energy}

\subsection*{Distances}
They match the paper:

\begin{align}
  d_{uv} &= d_{uv}^{xy} + d_{uv}^z,
          && \forall (u,v)\in\mathcal E^{ext}, \tag{DIST3}\\
  d_{uv}^{xy} &\ge \|(P_u^x-P_v^x,\; P_u^y-P_v^y)\|_2,
          && \forall (u,v)\in\mathcal E^{ext}, \tag{DIST4 (SOC)}\\
  d_{uv}^z &\ge |P_u^z - P_v^z|,
          && \forall (u,v)\in\mathcal E^{ext}. \tag{DIST5-6}
\end{align}

\subsection*{Inter-region energy}

The inter-region energy in the implementation follows the same structure
as the paper (with $\varphi_{rr'}$ and $\psi_{rr'}$), but all
bilinear products are explicitly linearized using the functions
\texttt{\_lin\_bb} (binary $\times$ binary) and \texttt{\_lin\_bc}
(binary $\times$ continuous).

The general implemented form is:
\begin{align}
  w_{uv} &= \underbrace{B^{xy}_{r1} \varphi_{rr'} d_{uv}^{xy} 
           \sum_{t\in\mathcal T_r} \alpha_t \rho_t}_{\text{linearized}}
           + \underbrace{B^z \varphi_{rr'}\psi_{rr'} d_{uv}^z
           \tfrac12(\rho_r+\rho_{r'})}_{\text{linearized}} \nonumber\\
         &\quad + \underbrace{B^{xy}_{r2} (1-\varphi_{rr'}) d_{uv}^{xy}
           \sum_{t\in\mathcal T_{r'}} \alpha_t \rho_t}_{\text{linearized}}
           + \underbrace{B^z (1-\varphi_{rr'})(1-\psi_{rr'}) d_{uv}^z
           \tfrac12(\rho_r+\rho_{r'})}_{\text{linearized}}, \tag{INTER-ENER}
\end{align}
where each bilinear term is replaced by an auxiliary variable with its
4 linearization constraints.

\section{Depot energy}

The implementation uses a simplified model with constant density $\rho_0$:

\begin{align}
  d_{v_0 v} &= d_{v}^{xy} + d_v^z,
            && \forall v\in\mathcal V', \tag{DEP-DIST}\\
  d_{v}^{xy} &\ge \|(P_v^x - d_x,\; P_v^y - d_y)\|_2,
            && \forall v\in\mathcal V', \tag{DEP-SOC}\\
  d_v^z &\ge |P_v^z|,
            && \forall v\in\mathcal V', \tag{DEP-Z}\\
  w_{v_0 v} &= E^{xy}\rho_0 \nu^{xy} d^{xy} + \tfrac12 E^z \rho_0 \nu^z d^z,
            && \forall v\in\mathcal V', \tag{DEP-ENER}\\
  w_{v v_0} &= E^{xy}\rho_0 \nu^{xy} d^{xy},
            && \forall v\in\mathcal V'. \tag{DEP-ENER-S}
\end{align}

\noindent \textbf{Difference:} In the paper, the ascent energy ($v_0\to v$)
considers the variable density $\rho_r + \rho_0$; the implementation uses only
$\rho_0$ as an approximation.

\section{Endurance constraint}

\begin{align}
  \sum_{u,v\in\mathcal V} \ell_{uv}^o &\le W^{max},
        && \forall o\in\mathcal O, \tag{END}\\
  \ell_{uv}^o &\le x_{uv}^o \cdot M,
        && \forall (u,v)\in\mathcal E,\; \forall o\in\mathcal O, \tag{END-LIN1}\\
  \ell_{uv}^o &\le w_{uv},
        && \forall (u,v)\in\mathcal E,\; \forall o\in\mathcal O, \tag{END-LIN2}\\
  \ell_{uv}^o &\ge w_{uv} - (1-x_{uv}^o) M,
        && \forall (u,v)\in\mathcal E,\; \forall o\in\mathcal O, \tag{END-LIN3}\\
  \ell_{uv}^o &\ge 0,
        && \forall (u,v)\in\mathcal E,\; \forall o\in\mathcal O. \tag{END-LIN4}
\end{align}
where $\ell_{uv}^o$ linearizes the product $x_{uv}^o \cdot w_{uv}$.

\section{Objective function}

\begin{equation}
  \min W = \sum_{u,v\in\mathcal V}\sum_{o\in\mathcal O}
  \ell_{uv}^o, \qquad \ell_{uv}^o \text{ linearizes } x_{uv}^o \cdot d_{uv}. \tag{OBI}
\end{equation}

\noindent \textbf{Difference:} The paper minimizes total distance;
the implementation does exactly the same, but with explicit
linearization of the product $x\cdot d$.

\section{Edge-based alternative (EdgesModel)}\label{sec:edge-model}

Alongside \texttt{RingsModel}, the codebase implements an alternative formulation
in \texttt{mip\_edges.py} (class \texttt{EdgesModel}) that replaces vertex-visitation
variables $y_v^o$ with edge-coverage variables $y_{uv}^o$ defined on the set of
intra-ring edges $\mathcal E^{int}$, following a perspective akin to the rural
postman problem.

\subsection*{Variables}

The binary variables $y_{uv}^o$ indicate whether the intra-ring edge $(u,v)$ is
covered in operation $o$. The integer variable $k^o$ counts the number of
covered edges instead of the number of visited vertices. All other variables
($x_{vv'}^o$, $\zeta^o$, $\alpha_t$, $\mu_v^{ts}$, $\gamma_v$, $\lambda_v$,
$P_v$, etc.) remain unchanged. The MTZ auxiliary variables $u_v^o$ are omitted.

\subsection*{Modified constraints}

The flow conservation (DP2) and depot constraints (DP1, DP3) are kept. The
degree constraints DP4--DP5 are relaxed to $\le 1$ (at most one edge leaves
or enters each vertex per operation). The vertex-visitation equalities
DP6 are removed, replaced by edge-coverage linking constraints:

\begin{align}
  \sum_{o\in\mathcal O} y_{uv}^o &= 1,
    &&\forall (u,v)\in\mathcal E^{int}, \tag{EC1}\\
  y_{uv}^o &= x_{uv}^o,
    &&\forall (u,v)\in\mathcal E^{int},\; \forall o\in\mathcal O. \tag{EC2}
\end{align}

\subsection*{Subtour elimination: DFJ lazy cuts}

The MTZ subtour elimination is not used. Instead, the model enforces subtour
elimination via the exponential family of Dantzig--Fulkerson--Johnson (DFJ)
constraints, separated dynamically through a lazy constraint callback:

\begin{equation}
  \sum_{u\in S,\,v\notin S} x_{uv}^o \ge 1,
    \qquad \forall S\subsetneq\mathcal V': v_0\notin S,\; \forall o\in\mathcal O,
    \tag{DFJ}
\end{equation}

At each feasible integer incumbent, a breadth-first search from the depot
identifies, for each operation $o$, the connected components of the subgraph
induced by edges with $x_{uv}^o = 1$. Any component that does not contain the
depot triggers the addition of cut~(DFJ). The callback is implemented using
Gurobi's \texttt{MIPSOL} callback with \texttt{LazyConstraints=1}.

\subsection*{Idle operations}

An additional constraint forces all traversal variables to zero in idle
operations ($\zeta^o = 1$), preventing the lazy callback from adding
spurious cuts:

\begin{equation}
  \sum_{u,v} x_{uv}^o \le |\mathcal V|^2\,(1 - \zeta^o),
    \qquad \forall o\in\mathcal O. \tag{IDLE}
\end{equation}

\subsection*{Relationship to the paper}

The alternative edge-based formulation described in Section~4.2 of the paper
matches this implementation exactly: the paper's relaxed degree constraints
(D$'_{\text{path}}$-C$_4$, D$'_{\text{path}}$-C$_5$) correspond to the
$\le 1$ degree bounds above; the DFJ lazy cut separation is described in
the paper under ``Drone path constraints (edge-based)''; and the edge-coverage
linking (EC$_1$, EC$_{2a}$, EC$_3$) is identical.

\section{Summary of differences}

\begin{center}
\begin{tabular}{lll}
\toprule
Aspect & Paper & Implementation \\
\midrule
Intra-RL distance & $\beta_{vv'}^{ss'} = \mu_v^s \mu_{v'}^{s'}$ (bilinear) & $dfs_v$ with $\eta_v^{ts} = \mu_v^{ts}\gamma_v$ linearized \\
Intra-RL energy & Analogous to distance & Same $efs_v$ strategy \\
Subtours (vertex-based) & MTZ ($u_v^o$) & MTZ ($u_v^o$) \\
Subtours (edge-based)   & DFJ (lazy cuts)   & DFJ (lazy cuts)   \\
START/END fixing & Implicit ($\lambda$) & Explicit (LC5b, LC8b) \\
Depot energy $\uparrow$ & $\rho_r + \rho_0$ variable & $\rho_0$ constant \\
Linearization & $\beta$ bilinear not linearized & All bilinearities explicitly linearized \\
Number of heights & Multiple heights with $\varphi,\psi$ & Multiple heights with $\varphi,\psi$ \\
\bottomrule
\end{tabular}
\end{center}

The main contribution of the implementation is the elimination of
bilinear products $\mu_v^s\mu_{v'}^{s'}$ through the $dfs_v$-based
formulation, which avoids $O(|\mathcal V'|^2|\mathcal S_t|^2)$ variables $\beta$.

\end{document}
